\documentclass[addpoints]{exam}
\usepackage{amsmath}
\usepackage{amssymb}
\usepackage{color}
\usepackage{siunitx}

% \printanswers

\newcommand{\event}{Trigonometry}
\newcommand{\tournament}{}
\def\type{0}

\setlength\fillinlinelength{\textwidth}
\CorrectChoiceEmphasis{\color{red}\bfseries}
\SolutionEmphasis{\color{red}\bfseries}

\setlength\answerclearance{2pt}
\pagestyle{headandfoot}
\runningheadrule
\runningheader{\event}{\tournament}{\ifnum\type=1 Class Set Test \else Full Name:\hspace{1cm} \fi}
\runningfootrule
\runningfooter{}{Page \thepage\ of \numpages}{}

\begin{document}
\begin{center}
    {\huge Grade 11 Foundations for College Math \\ \event
    \ifprintanswers \space Key
    \else \space \ifnum\type<2 Test \fi \ifnum\type=1 \textbf{(CLASS SET)} \fi
          \ifnum\type=2 Answer Sheet \fi
    \fi \\ \vspace{3mm}\tournament\vspace{1mm}}
\end{center}

\vspace{.25\textheight}

\begin{center}
    First Name: \ifprintanswers
    \underline{\hspace{3.5cm}}\textbf{KEY}\underline{\hspace{5.5cm}}\\\ \vspace{5mm}
    \else \underline{\hspace{10cm}}\\ \vspace{5mm} \fi
    Last Name: \underline{\hspace{10cm}}\\ \vspace{5mm}
\end{center}

\noindent\textbf{\underline{Directions:}}
\begin{itemize}
    \ifnum\type=1 \item \textbf{This test is a class set.} Do not write on this paper. \fi
    \item Show all work for full marks. Round dollar amounts to the nearest cent.
    \item The test is designed to be completed in 75 minutes.
\end{itemize}
\begin{center}
\textit{For grading use only}\\
\multirowgradetable{1}[pages]
\end{center}

\newpage
\begin{questions}

\section*{Part A --- Multiple Choice (15 marks)}
\vspace{5mm}

% Q1
\question[1]
In a right triangle with legs 3 and 4 and hypotenuse 5, $\sin\theta$ for the angle opposite the side of length 3 is:
\begin{choices}
    \choice $\dfrac{4}{5}$
    \correctchoice $\dfrac{3}{5}$
    \choice $\dfrac{3}{4}$
    \choice $\dfrac{5}{3}$
\end{choices}

% Q2
\question[1]
SOH-CAH-TOA defines $\tan\theta$ as:
\begin{choices}
    \choice $\dfrac{\text{hypotenuse}}{\text{adjacent}}$
    \choice $\dfrac{\text{adjacent}}{\text{opposite}}$
    \correctchoice $\dfrac{\text{opposite}}{\text{adjacent}}$
    \choice $\dfrac{\text{opposite}}{\text{hypotenuse}}$
\end{choices}

% Q3
\question[1]
In a right triangle, the side adjacent to a \ang{35} angle is \SI{10}{cm}. The hypotenuse is closest to:
\begin{choices}
    \choice \SI{8.19}{cm}
    \correctchoice \SI{12.21}{cm}
    \choice \SI{7.00}{cm}
    \choice \SI{14.28}{cm}
\end{choices}

% Q4
\question[1]
An angle of elevation is measured:
\begin{choices}
    \choice Below the horizontal.
    \correctchoice Above the horizontal.
    \choice From the vertical downward.
    \choice From the vertical upward.
\end{choices}

% Q5
\question[1]
A person standing \SI{50}{m} from the base of a tree looks up at an angle of elevation of \ang{40}. What is the approximate height of the tree?
\begin{choices}
    \choice \SI{32.14}{m}
    \correctchoice \SI{41.95}{m}
    \choice \SI{65.27}{m}
    \choice \SI{59.81}{m}
\end{choices}

% Q6
\question[1]
The Sine Law states:
\begin{choices}
    \choice $\dfrac{a}{\sin A} = \dfrac{b}{\cos B}$
    \correctchoice $\dfrac{a}{\sin A} = \dfrac{b}{\sin B} = \dfrac{c}{\sin C}$
    \choice $a^2 = b^2 + c^2 - 2bc\cos A$
    \choice $\sin A = \sin B + \sin C$
\end{choices}

% Q7
\question[1]
The Cosine Law is used when you have:
\begin{choices}
    \choice Two angles and one side (AAS).
    \choice Two angles and the included side (ASA).
    \correctchoice Two sides and the included angle (SAS), or three sides (SSS).
    \choice One angle and its opposite side.
\end{choices}

% Q8
\question[1]
In triangle $ABC$, $a = 7$, $b = 9$, $\angle A = 35°$. Using the Sine Law, $\sin B$ equals:
\begin{choices}
    \choice $\dfrac{7\sin 35°}{9}$
    \correctchoice $\dfrac{9\sin 35°}{7}$
    \choice $\dfrac{9}{7\sin 35°}$
    \choice $\dfrac{7}{9\sin 35°}$
\end{choices}

% Q9
\question[1]
In $\triangle PQR$, $p = 8$, $q = 6$, $r = 5$. To find $\angle P$ using the Cosine Law:
\begin{choices}
    \correctchoice $\cos P = \dfrac{q^2 + r^2 - p^2}{2qr}$
    \choice $\cos P = \dfrac{p^2 + r^2 - q^2}{2pr}$
    \choice $\cos P = \dfrac{p^2 + q^2 - r^2}{2pq}$
    \choice $\cos P = \dfrac{q^2 - r^2 - p^2}{2qr}$
\end{choices}

% Q10
\question[1]
Which of these is the Cosine Law formula for side $c$?
\begin{choices}
    \choice $c = a + b - 2ab\cos C$
    \correctchoice $c^2 = a^2 + b^2 - 2ab\cos C$
    \choice $c^2 = a^2 - b^2 + 2ab\cos C$
    \choice $c^2 = a^2 + b^2 + 2ab\cos C$
\end{choices}

% Q11
\question[1]
A ladder leans against a wall. The ladder is \SI{6}{m} long and makes a \ang{70} angle with the ground. How high up the wall does it reach?
\begin{choices}
    \choice \SI{2.05}{m}
    \correctchoice \SI{5.64}{m}
    \choice \SI{6.38}{m}
    \choice \SI{4.10}{m}
\end{choices}

% Q12
\question[1]
An angle of depression of \ang{25} is observed from the top of a cliff to a boat. If the cliff is \SI{80}{m} tall, the horizontal distance to the boat is approximately:
\begin{choices}
    \choice \SI{37.32}{m}
    \correctchoice \SI{171.56}{m}
    \choice \SI{88.28}{m}
    \choice \SI{190.01}{m}
\end{choices}

% Q13
\question[1]
In $\triangle ABC$, $\angle A = 50°$, $\angle B = 70°$, and $a = 10$. Find $b$ using the Sine Law. (Round to 1 decimal place.)
\begin{choices}
    \choice 7.7
    \correctchoice 12.3
    \choice 10.0
    \choice 9.4
\end{choices}

% Q14
\question[1]
$\triangle XYZ$ has $x = 5$, $y = 7$, $\angle Z = 60°$. Using the Cosine Law, $z^2$ equals:
\begin{choices}
    \choice $25 + 49 + 35$
    \choice $25 + 49$
    \correctchoice $25 + 49 - 35$
    \choice $25 - 49 + 35$
\end{choices}

% Q15
\question[1]
In which situation must the Cosine Law be used (Sine Law alone is insufficient)?
\begin{choices}
    \choice AAS — two angles and a non-included side.
    \choice ASA — two angles and the included side.
    \correctchoice SSS — all three sides known, find an angle.
    \choice SAA — one side and two non-adjacent angles.
\end{choices}

\section*{Part B --- Short Answer (33 marks)}

% Q16
\question[6]
\textbf{Right Triangle Trigonometry.} Solve each of the following. Round lengths to one decimal place and angles to the nearest degree.

\begin{parts}
    \part[2] In right $\triangle ABC$ ($\angle C = 90°$), $\angle A = 38°$ and $c = 15$ cm. Find side $a$ (opposite $\angle A$).
    \begin{solution}
        $\sin A = \dfrac{a}{c} \Rightarrow a = c\sin A = 15\sin 38°$\\
        $a = 15 \times 0.6157 \approx \SI{9.2}{cm}$
    \end{solution}

    \part[2] In right $\triangle DEF$ ($\angle F = 90°$), $d = 8$ cm (side opposite $\angle D$) and $e = 11$ cm (side opposite $\angle E$). Find $\angle D$.
    \begin{solution}
        $\tan D = \dfrac{\text{opposite}}{\text{adjacent}} = \dfrac{d}{e} = \dfrac{8}{11}$\\
        $\angle D = \arctan\!\left(\dfrac{8}{11}\right) \approx \arctan(0.7273) \approx 36°$
    \end{solution}

    \part[2] In right $\triangle PQR$ ($\angle R = 90°$), $\angle Q = 55°$ and $p = 20$ cm (the leg adjacent to $\angle Q$). Find the hypotenuse $q$.
    \begin{solution}
        $\cos Q = \dfrac{p}{q} \Rightarrow q = \dfrac{p}{\cos Q} = \dfrac{20}{\cos 55°}$\\
        $q = \dfrac{20}{0.5736} \approx \SI{34.9}{cm}$
    \end{solution}
\end{parts}

% Q17
\question[5]
\textbf{Angles of Elevation and Depression.}

\begin{parts}
    \part[2] A flagpole casts a shadow \SI{12}{m} long when the angle of elevation of the sun is \ang{48}. How tall is the flagpole? Round to one decimal place.
    \begin{solution}
        $\tan 48° = \dfrac{h}{12}$\\
        $h = 12\tan 48° = 12 \times 1.1106 \approx \SI{13.3}{m}$
    \end{solution}

    \part[3] An airplane is flying at an altitude of \SI{3500}{m}. From a point on the ground directly below the plane's previous position, the angle of elevation to the plane is \ang{22}. How far has the plane flown horizontally from that point? Round to the nearest metre.
    \begin{solution}
        $\tan 22° = \dfrac{3500}{d}$\\
        $d = \dfrac{3500}{\tan 22°} = \dfrac{3500}{0.4040} \approx \SI{8663}{m}$
    \end{solution}
\end{parts}

% Q18
\question[6]
\textbf{The Sine Law.}

\begin{parts}
    \part[3] In $\triangle ABC$, $\angle A = 42°$, $\angle B = 73°$, and $a = 14$ cm. Find side $b$. Round to one decimal place.
    \begin{solution}
        $\dfrac{b}{\sin B} = \dfrac{a}{\sin A}$\\
        $b = \dfrac{a \sin B}{\sin A} = \dfrac{14 \sin 73°}{\sin 42°}$\\
        $b = \dfrac{14 \times 0.9563}{0.6691} \approx \dfrac{13.388}{0.6691} \approx \SI{20.0}{cm}$
    \end{solution}

    \part[3] In $\triangle PQR$, $p = 10$ cm, $r = 7$ cm, and $\angle R = 30°$. Find $\angle P$ to the nearest degree.
    \begin{solution}
        $\dfrac{\sin P}{p} = \dfrac{\sin R}{r}$\\
        $\sin P = \dfrac{p \sin R}{r} = \dfrac{10 \sin 30°}{7} = \dfrac{10 \times 0.5}{7} = \dfrac{5}{7} \approx 0.7143$\\
        $\angle P = \arcsin(0.7143) \approx 46°$
    \end{solution}
\end{parts}

% Q19
\question[8]
\textbf{The Cosine Law.}

\begin{parts}
    \part[4] In $\triangle DEF$, $d = 9$ cm, $e = 12$ cm, and $\angle F = 65°$. Find side $f$. Round to one decimal place.
    \begin{solution}
        $f^2 = d^2 + e^2 - 2de\cos F$\\
        $f^2 = 81 + 144 - 2(9)(12)\cos 65°$\\
        $f^2 = 225 - 216 \times 0.4226$\\
        $f^2 = 225 - 91.28 = 133.72$\\
        $f = \sqrt{133.72} \approx \SI{11.6}{cm}$
    \end{solution}

    \part[4] In $\triangle XYZ$, $x = 6$ cm, $y = 8$ cm, $z = 10$ cm. Find $\angle Z$ to the nearest degree.
    \begin{solution}
        $\cos Z = \dfrac{x^2 + y^2 - z^2}{2xy} = \dfrac{36 + 64 - 100}{2(6)(8)} = \dfrac{0}{96} = 0$\\
        $\angle Z = \arccos(0) = 90°$\\
        (This is a right triangle with hypotenuse $z = 10$.)
    \end{solution}
\end{parts}

% Q20
\question[8]
\textbf{Practical Application: Width of a Lake.} A surveyor wants to find the width $AB$ of a lake. She stands at point $C$ on one shore and measures $CA = \SI{85}{m}$ and $CB = \SI{110}{m}$. She also measures the angle between the two lines: $\angle ACB = 52°$.

\begin{parts}
    \part[4] Use the Cosine Law to find the width $AB$ of the lake. Round to the nearest metre.
    \begin{solution}
        Let $a = CB = 110$, $b = CA = 85$, $\angle C = 52°$, and $c = AB$.\\
        $c^2 = a^2 + b^2 - 2ab\cos C$\\
        $c^2 = 110^2 + 85^2 - 2(110)(85)\cos 52°$\\
        $c^2 = 12100 + 7225 - 18700 \times 0.6157$\\
        $c^2 = 19325 - 11513.59 = 7811.41$\\
        $c = \sqrt{7811.41} \approx \SI{88}{m}$\\
        The width of the lake is approximately \SI{88}{m}.
    \end{solution}

    \part[2] The surveyor also wants to find $\angle CAB$. Use the Sine Law to find this angle.
    \begin{solution}
        $\dfrac{\sin(\angle CAB)}{CB} = \dfrac{\sin C}{AB}$\\
        $\sin(\angle CAB) = \dfrac{CB \cdot \sin C}{AB} = \dfrac{110 \sin 52°}{88}$\\
        $\sin(\angle CAB) = \dfrac{110 \times 0.7880}{88} \approx \dfrac{86.68}{88} \approx 0.9850$\\
        $\angle CAB = \arcsin(0.9850) \approx 80°$
    \end{solution}

    \part[2] Using the angles of the triangle, find $\angle ABC$.
    \begin{solution}
        Sum of angles in a triangle $= 180°$.\\
        $\angle ABC = 180° - 52° - 80° = 48°$
    \end{solution}
\end{parts}

\end{questions}
\end{document}
